\section{Diseño e implementación del parser}

\subsection{Decisión de diseño: descenso recursivo}

El enunciado permite implementar el analizador predictivo mediante descenso recursivo o mediante una tabla LL(1). Elegimos \textbf{descenso recursivo}, porque el código queda casi idéntico a la gramática: \textbf{cada no terminal es un método} de la clase \texttt{Parser}, y cada método decide qué producción aplicar mirando solo el token actual (un token de anticipación), exactamente como indica la tabla de la sección 6.

La tabla LL(1) no se usa dentro del parser, pero tampoco queda sin validar: el módulo \texttt{gramatica.py} recalcula FIRST, FOLLOW y la tabla de forma automática a partir de las 27 producciones, y confirma que coinciden con los cálculos hechos a mano en las secciones 5 y 6 (el comando \texttt{python -m gamerules.gramatica} informa que no hay conflictos). El parser no depende de ninguna biblioteca que genere el análisis: está escrito completamente a mano.

\subsection{Organización del código}

\begin{table}[H]
\centering
\caption{Módulos del proyecto.}
\small
\renewcommand{\arraystretch}{1.3}
\begin{tabularx}{\textwidth}{|l|X|}
\hline
\textbf{Archivo} & \textbf{Responsabilidad} \\ \hline
\texttt{main.py} & Programa principal: línea de comandos, lectura de archivos y modo interactivo \\ \hline
\texttt{gamerules/lexico.py} & Analizador léxico: convierte el texto en una lista de tokens con línea y columna \\ \hline
\texttt{gamerules/sintactico.py} & Parser LL(1) por descenso recursivo, con detección y recuperación de errores \\ \hline
\texttt{gamerules/arbol.py} & Clase \texttt{Nodo}, impresión del árbol en texto y exportación a Graphviz \\ \hline
\texttt{gamerules/gramatica.py} & Verificación automática de FIRST, FOLLOW y la tabla LL(1) (no la usa el parser) \\ \hline
\texttt{pruebas/} & 30 casos de prueba y el programa que los ejecuta \\ \hline
\end{tabularx}
\end{table}

El flujo completo es: texto de entrada $\rightarrow$ \texttt{analizador\_lexico} $\rightarrow$ lista de tokens $\rightarrow$ \texttt{Parser} $\rightarrow$ árbol sintáctico o lista de errores. La función \texttt{analizar(texto)} encadena las dos fases y retorna \texttt{(arbol, errores)}: si la entrada es válida el árbol es un objeto \texttt{Nodo} y la lista de errores está vacía; si no, el árbol es \texttt{None}.

\subsection{Cambios al analizador léxico}

Partimos del analizador léxico entregado por el docente y le hicimos dos modificaciones:

\begin{enumerate}
    \item Agregamos el token \texttt{COMA}, necesario para separar los parámetros de una acción.
    \item Cada token ahora es una tupla con nombre \texttt{Token(tipo, lexema, linea, columna)}. Para llevar la cuenta de la línea, el salto de línea pasó a ser un token propio (\texttt{SALTO}) que el léxico consume sin entregarlo al parser. Gracias a esto el parser puede indicar \emph{dónde} ocurre cada error.
\end{enumerate}

Además, el error léxico dejó de ser un \texttt{ValueError} genérico: ahora se lanza una excepción \texttt{ErrorLexico} que incluye la posición del símbolo no reconocido. Al final de la lista de tokens, el parser agrega el token especial \texttt{\$} que marca el fin de la entrada.

\subsection{Correspondencia entre gramática y código}

La siguiente tabla relaciona cada no terminal con su método y con el conjunto que el método consulta para decidir. Estos conjuntos están definidos como constantes al inicio de \texttt{sintactico.py} (por ejemplo \texttt{FOLLOW\_TERMINO\_PRIMA}) y son los mismos que se calcularon en la sección 5.

\begin{table}[H]
\centering
\caption{Relación entre producciones, métodos del parser y conjuntos de decisión.}
\small
\renewcommand{\arraystretch}{1.3}
\begin{tabularx}{\textwidth}{|l|l|c|X|}
\hline
\textbf{No terminal} & \textbf{Método} & \textbf{Prod.} & \textbf{Cómo decide} \\ \hline
\texttt{<programa>} & \texttt{programa} & 1 & Una sola alternativa: una regla y luego \texttt{<programa'>} \\ \hline
\texttt{<programa'>} & \texttt{programa\_prima} & 2, 3 & \texttt{WHEN} $\rightarrow$ 2; \texttt{\$} $\rightarrow$ 3 \\ \hline
\texttt{<regla>} & \texttt{regla} & 4 & Consume \texttt{WHEN}, condición, \texttt{THEN}, acción y \texttt{;} \\ \hline
\texttt{<condicion>} & \texttt{condicion} & 5 & Una sola alternativa \\ \hline
\texttt{<condicion'>} & \texttt{condicion\_prima} & 6, 7 & \texttt{OR} $\rightarrow$ 6; FOLLOW $=\{\texttt{THEN}, \texttt{)}\}$ $\rightarrow$ 7 \\ \hline
\texttt{<termino>} & \texttt{termino} & 8 & Una sola alternativa \\ \hline
\texttt{<termino'>} & \texttt{termino\_prima} & 9, 10 & \texttt{AND} $\rightarrow$ 9; FOLLOW $=\{\texttt{OR}, \texttt{THEN}, \texttt{)}\}$ $\rightarrow$ 10 \\ \hline
\texttt{<factor>} & \texttt{factor} & 11--13 & \texttt{NOT} $\rightarrow$ 11; \texttt{(} $\rightarrow$ 12; \texttt{ID} $\rightarrow$ 13 \\ \hline
\texttt{<comparacion>} & \texttt{comparacion} & 14, 15 & \texttt{OP\_REL} $\rightarrow$ 14; FOLLOW $=\{\texttt{AND}, \texttt{OR}, \texttt{THEN}, \texttt{)}\}$ $\rightarrow$ 15 \\ \hline
\texttt{<valor>} & \texttt{valor} & 16--19 & \texttt{NUM}, \texttt{TRUE}, \texttt{FALSE} o \texttt{ID} \\ \hline
\texttt{<accion>} & \texttt{accion} & 20 & Una sola alternativa: exige un \texttt{ID} como verbo \\ \hline
\texttt{<objeto>} & \texttt{objeto} & 21, 22 & \texttt{ID} $\rightarrow$ 21; FOLLOW $=\{\texttt{(}, \texttt{;}\}$ $\rightarrow$ 22 \\ \hline
\texttt{<parametros>} & \texttt{parametros} & 23, 24 & \texttt{(} $\rightarrow$ 23; \texttt{;} $\rightarrow$ 24 \\ \hline
\texttt{<argumentos>} & \texttt{argumentos} & 25 & Exige que el token esté en FIRST(\texttt{<valor>}) \\ \hline
\texttt{<argumentos'>} & \texttt{argumentos\_prima} & 26, 27 & \texttt{,} $\rightarrow$ 26; \texttt{)} $\rightarrow$ 27 \\ \hline
\end{tabularx}
\end{table}

Las producciones $\varepsilon$ se aplican cuando el token actual pertenece al FOLLOW del no terminal; si no pertenece a ninguno de los dos casos, la celda correspondiente de la tabla está vacía y el método informa un error sintáctico (sección 9).

\subsection{Ejemplos del código}

El método de \texttt{<condicion'>} implementa las producciones 6 y 7. Se ve cómo la alternativa $\varepsilon$ se elige cuando el token pertenece a FOLLOW, y cualquier otro token es un error:

\begin{lstlisting}[language=Python]
def condicion_prima(self):
    nodo = Nodo("<condicion'>")
    if self.actual.tipo == 'OR':
        # 6: <condicion'> -> OR <termino> <condicion'>
        nodo.agregar(self.coincidir('OR'))
        nodo.agregar(self.termino())
        nodo.agregar(self.condicion_prima())
    elif self.actual.tipo in FOLLOW_CONDICION_PRIMA:
        # 7: <condicion'> -> epsilon
        nodo.agregar(Nodo('eps'))
    else:
        raise self.error(self.esperado_fin_condicion())
    return nodo
\end{lstlisting}

El método de \texttt{<factor>} implementa las tres alternativas de las producciones 11, 12 y 13. Como sus FIRST son disjuntos (\texttt{NOT}, \texttt{(} e \texttt{ID}), basta un token para elegir:

\begin{lstlisting}[language=Python]
def factor(self):
    nodo = Nodo('<factor>')
    if self.actual.tipo == 'NOT':              # 11
        nodo.agregar(self.coincidir('NOT'))
        nodo.agregar(self.factor())
    elif self.actual.tipo == 'PARENT_IZQ':     # 12
        nodo.agregar(self.coincidir('PARENT_IZQ'))
        nodo.agregar(self.condicion())
        nodo.agregar(self.coincidir('PARENT_DER', ...))
    elif self.actual.tipo == 'ID':             # 13
        nodo.agregar(self.coincidir('ID'))
        nodo.agregar(self.comparacion())
    else:
        raise self.error("una condicion (un identificador, NOT o '(')")
    return nodo
\end{lstlisting}

\subsection{Funciones auxiliares}

\begin{itemize}
    \item \texttt{coincidir(tipo)}: si el token actual es del tipo esperado lo consume y crea una hoja del árbol; si no, lanza un \texttt{ErrorSintactico}. Es la operación \emph{match} de los analizadores predictivos.
    \item \texttt{avanzar()}: pasa al siguiente token (sin pasar nunca del marcador \texttt{\$}).
    \item \texttt{error(esperado)}: construye el mensaje de error con la posición (sección 9).
    \item \texttt{sincronizar()}: recuperación en modo pánico después de un error (sección 9).
\end{itemize}

\subsection{Ambigüedad y precedencia en el código}

La estratificación de la gramática (sección 4) se refleja directamente en la estructura de las llamadas: \texttt{condicion} llama a \texttt{termino} y este llama a \texttt{factor}. Por eso \texttt{a OR b AND c} se analiza siempre como \texttt{a OR (b AND c)}: el \texttt{AND} se consume dentro de \texttt{termino}, antes de que \texttt{condicion\_prima} vea el \texttt{OR}. No hace falta ninguna regla de precedencia adicional en el código.

\newpage
